\section{Reproducibility Details}
\label{app:repro}

% OWNER: Jackson Wilke (methods/infra per proposal)

\paragraph{Pipeline.} The LSTM pipeline is
\texttt{data\_utils} $\to$ \texttt{train} $\to$ checkpoint $\to$
\texttt{evaluate} $\to$ \texttt{results/eval\_results.csv} $\to$ \texttt{plots}.
\texttt{data\_utils} builds the shared vocabulary and the fixed 20K held-out
indices; \texttt{minimal\_pairs} and \texttt{verb\_flip} construct the leak-free
test pairs from the clean corpus. The RNNG pipeline
(\texttt{run\_rnng\_datahub.py}) parses the clean corpus once with the Berkeley
Neural Parser, swaps the verb terminal per noise condition, runs the
\texttt{rnng-pytorch} preprocessing and training, and scores the shared minimal
pairs by beam search.

\paragraph{Cross-model comparison (H3).} The LSTM scores its 19{,}819 in-vocab
pairs via \texttt{src/evaluate.py}; the full-scale RNNG run scores 19{,}983
held-out pairs by beam search within \texttt{run\_rnng\_datahub.py}. We compare
overall accuracy across these near-identical sets directly and take the
within-model attractor gap as the primary cross-model contrast, since it is a
rate rather than a count and is robust to the small difference in covered pairs
(Section~\ref{sec:metrics}). \texttt{src/combined\_plots.py} reads both eval
CSVs and produces the side-by-side accuracy and attractor-gap figures.

\paragraph{Validity checks.} The RNNG run also emits a per-epoch validation
perplexity log (\texttt{results/rnng\_train\_ppl.csv}) to confirm convergence,
and a parser-quality summary plus a sample of parsed trees
(\texttt{results/rnng\_parse\_quality.json}, \texttt{rnng\_tree\_sample.txt}) for
hand-inspection, since parser quality on lowercased, anonymized tokens bounds the
RNNG's ceiling.

\begin{table}[t]
  \centering
  \small
  \begin{tabular}{lcc}
    \toprule
    Hyperparameter & LSTM & RNNG \\
    \midrule
    Embedding / word dim   & 256 & 256 \\
    Hidden dim             & 256 & 256 \\
    Layers                 & 2   & 2 \\
    Dropout                & 0.2 & 0.3 \\
    Tied embeddings        & yes & --- \\
    Optimizer              & Adam & Adam \\
    Learning rate          & $10^{-3}$ & $10^{-3}$ \\
    Batch size             & 64  & 128 (15K-token cap) \\
    Gradient clip          & 1.0 & --- \\
    Epochs                 & 5   & 8 \\
    Sentences / condition  & 150K & 150K \\
    Vocabulary             & 10K (shared) & 10K (shared) \\
    Decoding               & next-token & beam 100 / word-beam 20 \\
    \bottomrule
  \end{tabular}
  \caption{Hyperparameters for the LSTM and RNNG. Embedding, hidden size,
    optimizer, and training-set size are matched; the RNNG additionally consumes
    benepar parse trees and is decoded by word-synchronous beam search.}
  \label{tab:hparams}
\end{table}

\section{Additional Figures}
\label{app:figures}

Figure~\ref{fig:gap_trajectory} shows the attractor-gap trajectory over training
(companion to the accuracy trajectory in Section~\ref{sec:results}): the gap
shrinks over training in low-noise conditions as the hierarchical rule
consolidates, but grows in the 50\% condition as additional exposure entrenches
the linear heuristic.

\begin{figure}[t]
  \centering
  \includegraphics[width=\linewidth]{gap_trajectory.png}
  \caption{Attractor gap over training steps, one curve per noise condition.}
  \label{fig:gap_trajectory}
\end{figure}
